% Feedback:
% 1. Years of Experience — maybe consider two pages.
%    Write as much as you need to. Thanks you.
% 2. The first page is therefore prime real estate.
% 3. Skills section shorter.
% 4. Don't bother with portfolio. You have years of experience.
% 5.

%%%%%%%%%%%%%%%%%%%%%%%%%%%%%%%%%%%%%%%%%%%%%%%%%%%%%%%%%%%%%%%%%%%%%%%%
%%%%%%%%%%%%%%%%%%%%%% Simple LaTeX CV Template %%%%%%%%%%%%%%%%%%%%%%%%
%%%%%%%%%%%%%%%%%%%%%%%%%%%%%%%%%%%%%%%%%%%%%%%%%%%%%%%%%%%%%%%%%%%%%%%%

%%%%%%%%%%%%%%%%%%%%%%%%%%%%%%%%%%%%%%%%%%%%%%%%%%%%%%%%%%%%%%%%%%%%%%%%
%% NOTE: If you find that it says                                     %%
%%                                                                    %%
%%                           1 of ??                                  %%
%%                                                                    %%
%% at the bottom of your first page, this means that the AUX file     %%
%% was not available when you ran LaTeX on this source. Simply RERUN  %%
%% LaTeX to get the ``??'' replaced with the number of the last page  %%
%% of the document. The AUX file will be generated on the first run   %%
%% of LaTeX and used on the second run to fill in all of the          %%
%% references.                                                        %%
%%%%%%%%%%%%%%%%%%%%%%%%%%%%%%%%%%%%%%%%%%%%%%%%%%%%%%%%%%%%%%%%%%%%%%%%

%%%%%%%%%%%%%%%%%%%%%%%%%%%% Document Setup %%%%%%%%%%%%%%%%%%%%%%%%%%%%

% Don't like 10pt? Try 11pt or 12pt
\documentclass[10pt]{article}

% The automated optical recognition software used to digitize resume
% information works best with fonts that do not have serifs. This
% command uses a sans serif font throughout. Uncomment both lines (or at
% least the second) to restore a Roman font (i.e., a font with serifs).
%\usepackage{times}
%\renewcommand{\familydefault}{\sfdefault}

% This is a helpful package that puts math inside length specifications
\usepackage{calc}
\usepackage{comment}
\usepackage{graphicx}
\usepackage{stackengine}[2013-10-15]
\newcommand\textss[1]{\stackengine{.9ex}{}{\scriptsize#1}{O}{l}{F}{F}{L}}
% \usepackage{fontawesome}

% Simpler bibsection for CV sections
% (thanks to natbib for inspiration)
\makeatletter
\newlength{\bibhang}
\setlength{\bibhang}{1em} %1em}
\newlength{\bibsep}
 {\@listi \global\bibsep\itemsep \global\advance\bibsep by\parsep}
\newenvironment{bibsection}%
		{\begin{enumerate}{}{%
%        {\begin{list}{}{%
	   \setlength{\leftmargin}{\bibhang}%
	   \setlength{\itemindent}{-\leftmargin}%
	   \setlength{\itemsep}{\bibsep}%
	   \setlength{\parsep}{\z@}%
		\setlength{\partopsep}{0pt}%
		\setlength{\topsep}{0pt}}}
		{\end{enumerate}\vspace{-.6\baselineskip}}
%        {\end{list}\vspace{-.6\baselineskip}}
\makeatother

% Layout: Puts the section titles on left side of page
\reversemarginpar

%
%         PAPER SIZE, PAGE NUMBER, AND DOCUMENT LAYOUT NOTES:
%
% The next \usepackage line changes the layout for CV style section
% headings as marginal notes. It also sets up the paper size as either
% letter or A4. By default, letter was used. If A4 paper is desired,
% comment out the letterpaper lines and uncomment the a4paper lines.
%
% As you can see, the margin widths and section title widths can be
% easily adjusted.
%
% ALSO: Notice that the includefoot option can be commented OUT in order
% to put the PAGE NUMBER *IN* the bottom margin. This will make the
% effective text area larger.
%
% IF YOU WISH TO REMOVE THE ``of LASTPAGE'' next to each page number,
% see the note about the +LP and -LP lines below. Comment out the +LP
% and uncomment the -LP.
%
% IF YOU WISH TO REMOVE PAGE NUMBERS, be sure that the includefoot line
% is uncommented and ALSO uncomment the \pagestyle{empty} a few lines
% below.
%

%% Use these lines for letter-sized paper
%\usepackage[paper=letterpaper,
%            %includefoot, % Uncomment to put page number above margin
%            marginparwidth=1.2in,     % Length of section titles
%            marginparsep=.05in,       % Space between titles and text
%            margin=1in,               % 1 inch margins
%            includemp]{geometry}

% Use these lines for A4-sized paper
\usepackage[paper=a4paper,
			%includefoot, % Uncomment to put page number above margin
			marginparwidth=0mm,    % Length of section titles
			marginparsep=0mm,       % Space between titles and text
			margin=18mm,              % 25mm margins
			includemp]{geometry}

%% More layout: Get rid of indenting throughout entire document
\setlength{\parindent}{0in}

\usepackage[shortlabels]{enumitem}

%% Reference the last page in the page number
%
% NOTE: comment the +LP line and uncomment the -LP line to have page
%       numbers without the ``of ##'' last page reference)
%
% NOTE: uncomment the \pagestyle{empty} line to get rid of all page
%       numbers (make sure includefoot is commented out above)
%
\usepackage{fancyhdr,lastpage}
\pagestyle{fancy}
%\pagestyle{empty}      % Uncomment this to get rid of page numbers
\fancyhf{}\renewcommand{\headrulewidth}{0pt}
\fancyfootoffset{\marginparsep+\marginparwidth}
\newlength{\footpageshift}
\setlength{\footpageshift}
		  {0.5\textwidth+0.5\marginparsep+0.5\marginparwidth-2in}
\lfoot{\hspace{\footpageshift}%
	   \parbox{4in}{\, \hfill %
					\arabic{page} of \protect\pageref*{LastPage} % +LP
%                    \arabic{page}                               % -LP
					\hfill \,}}

% Finally, give us PDF bookmarks
\usepackage{color,hyperref}
\definecolor{darkblue}{rgb}{0.0,0.0,0.3}
\hypersetup{colorlinks,breaklinks,
			linkcolor=darkblue,urlcolor=darkblue,
			anchorcolor=darkblue,citecolor=darkblue}

%%%%%%%%%%%%%%%%%%%%%%%% End Document Setup %%%%%%%%%%%%%%%%%%%%%%%%%%%%


%%%%%%%%%%%%%%%%%%%%%%%%%%% Helper Commands %%%%%%%%%%%%%%%%%%%%%%%%%%%%

% The title (name) with a horizontal rule under it
% (optional argument typesets an object right-justified across from name
%  as well)
%
% Usage: \makeheading{name}
%        OR
%        \makeheading[right_object]{name}
%
% Place at top of document. It should be the first thing.
% If ``right_object'' is provided in the square-braced optional
% argument, it will be right justified on the same line as ``name'' at
% the top of the CV. For example:
%
%       \makeheading[\emph{Curriculum vitae}]{Your Name}
%
% will put an emphasized ``Curriculum vitae'' at the top of the document
% as a title. Likewise, a picture could be included:
%
%   \makeheading[\includegraphics[height=1.5in]{my_picutre}]{Your Name}
%
% the picture will be flush right across from the name.
\newcommand{\makeheading}[8][]{
	\hspace*{-\marginparsep minus \marginparwidth}%
	\begin{minipage}[t]{\textwidth+\marginparwidth+\marginparsep}%
		\begin{tabular}[t]{@{}l}
		{\hspace{-8pt} \Large \bfseries#2}\\
		\\
		{\normalsize \emph{#3}}
		\end{tabular}
		\hfill
		\begin{tabular}[t]{|l}
		#4  \\
		#5  \\
		#6 \\

		\end{tabular}
		\vspace{1pt}
		\\[-0.2\baselineskip]%

		\rule{\columnwidth}{1pt}%
	\end{minipage}
}

% The section headings
%
% Usage: \section{section name}
% \renewcommand{\section}[1]{\pagebreak[3]%
%     \hyphenpenalty=10000%
%     \vspace{1.8\baselineskip}%
%     \phantomsection\addcontentsline{toc}{section}{#1}%
%     \noindent\llap{\scshape\smash{\parbox[t]{\marginparwidth}{\raggedright #1}}}%
%     \vspace{-\baselineskip}\par}

\renewcommand{\section}[1]{
	\pagebreak[3]%
	\hyphenpenalty=10000%
	\vspace{\baselineskip}%
	\phantomsection\addcontentsline{toc}{section}{#1}%
	\noindent{\scshape #1}%
	\vspace{\baselineskip}\par
}

\newcommand{\sectioncont}[1]{
	\pagebreak[3]%
	\hyphenpenalty=10000%
	\vspace{\baselineskip}%
	\noindent{\scshape #1} (continuación)%
	\vspace{\baselineskip}\par
}

% An itemize-style list with lots of space between items
\newenvironment{outerlist}[1][\enskip\textbullet]{
\begin{itemize}[#1,leftmargin=*]}{\end{itemize}%
\vspace{-.6\baselineskip}
}

% An environment IDENTICAL to outerlist that has better pre-list spacing
% when used as the first thing in a \section
\newenvironment{lonelist}[1][\enskip\textbullet]{
	\begin{list}{#1}{%
	\setlength{\partopsep}{0pt}%
	\setlength{\topsep}{0pt}}
}{\end{list}\vspace{-.6\baselineskip}}

% An itemize-style list with little space between items
\newenvironment{innerlist}[1][\enskip\textbullet]{
	\begin{itemize}[#1,leftmargin=*,rightmargin=30pt,parsep=0pt,itemsep=1pt,topsep=2pt,partopsep=0pt]
}{\end{itemize}}

\newenvironment{innerlist2}[1][\enskip\textasteriskcentered]{
	\begin{itemize}[#1,leftmargin=*,parsep=0pt,itemsep=1pt,topsep=2pt,partopsep=0pt]
}{\end{itemize}}

% An environment IDENTICAL to innerlist that has better pre-list spacing
% when used as the first thing in a \section
\newenvironment{loneinnerlist}[1][\enskip\textbullet]%
		{\begin{itemize}[#1,leftmargin=*,parsep=0pt,itemsep=0pt,topsep=0pt,partopsep=0pt]}
		{\end{itemize}\vspace{-.6\baselineskip}}

% To add some paragraph space between lines.
% This also tells LaTeX to preferably break a page on one of these gaps
% if there is a needed pagebreak nearby.
\newcommand{\blankline}{\quad\pagebreak[3]}
\newcommand{\halfblankline}{\quad\vspace{-0.5\baselineskip}\pagebreak[3]}

% Uses hyperref to link DOI
\newcommand\doilink[1]{\href{http://dx.doi.org/#1}{#1}}
\newcommand\doi[1]{doi:\doilink{#1}}

% For \url{SOME_URL}, links SOME_URL to the url SOME_URL
\providecommand*\url[1]{\href{#1}{#1}}
% Same as above, but pretty-prints SOME_URL in teletype fixed-width font
\renewcommand*\url[1]{\href{#1}{\texttt{#1}}}

% For \email{ADDRESS}, links ADDRESS to the url mailto:ADDRESS
\providecommand*\email[1]{\href{mailto:#1}{#1}}
% Same as above, but pretty-prints ADDRESS in teletype fixed-width font
%\renewcommand*\email[1]{\href{mailto:#1}{\texttt{#1}}}

\newcommand{\nohref}[2]{#2}

%\providecommand\BibTeX{{\rm B\kern-.05em{\sc i\kern-.025em b}\kern-.08em
%    T\kern-.1667em\lower.7ex\hbox{E}\kern-.125emX}}
%\providecommand\BibTeX{{\rm B\kern-.05em{\sc i\kern-.025em b}\kern-.08em
%    \TeX}}
\providecommand\BibTeX{{B\kern-.05em{\sc i\kern-.025em b}\kern-.08em
	\TeX}}
\providecommand\Matlab{\textsc{MATLAB}}

% Work experience subsection with PDF navigation
% Usage: \worksubsection{Job Title}{Company (with links)}{Date Range}
\newcommand{\worksubsection}[3]{
	\phantomsection
	\addcontentsline{toc}{subsection}{#1 — #2}
	\setlength\tabcolsep{0pt}
	\noindent\textbf{#1} — #2 \hfill \emph{#3} \\
	\vspace{-6pt}
}

% Education subsection with PDF navigation
% Usage: \edusubsection{Degree}{Institution}{Honors/Grade (optional)}
\newcommand{\edusubsection}[3]{
	\phantomsection
	\addcontentsline{toc}{subsection}{#1 — #2}
	\setlength\tabcolsep{0pt}
	\noindent\textbf{#1} — \textbf{#2}%
	\if\relax\detokenize{#3}\relax\else\ (\emph{#3})\fi\\
	\vspace{-2pt}
}

%%%%%%%%%%%%%%%%%%%%%%%% End Helper Commands %%%%%%%%%%%%%%%%%%%%%%%%%%%

%%%%%%%%%%%%%%%%%%%%%%%%% Begin CV Document %%%%%%%%%%%%%%%%%%%%%%%%%%%%
\newlength{\rcollength}\setlength{\rcollength}{1.2in}%


\begin{document}
\makeheading{
		\href{https://joshuaprettyman.com/}{Joshua Prettyman, Ph.D.}
}
{Científico de Datos, Doctor en Matemáticas, Desarrollador Python}
% {Python Developer, Data Scientist, Mathematics Ph.D.}
{\email{joshua@prettyman.me}}
% {\href{https://github.com/DrPrettyman}{github.com/DrPrettyman}}
{\href{https://www.linkedin.com/in/joshuaprettyman/}{linkedin.com/in/joshuaprettyman}}
% {{\href{tel:+44(0)7470885347}{tel: 07470 885 347}}}
% {{\href{tel:+34711042515}{(+34) 711 042 151}}}
{\href{https://joshuaprettyman.com/}{joshuaprettyman.com}}



%\section{Objective}

%Insert text here if you want to
%\begin{innerlist}
%\item More information and auxiliary documents can be found at\\\url{http://www.tedpavlic.com/facjobsearch/}
%\end{innerlist}



\section{Resumen Profesional}
Desarrollé una plataforma SaaS impulsada por ML que aumentó la productividad en SEO en $20\times$, transformando las herramientas internas de Blink SEO de hojas de cálculo a un producto comercializable ahora utilizado por múltiples agencias: \href{https://macalytics.app}{\textit{Macaroni Software}}. Contratado como su primer Científico de Datos, construí la plataforma full-stack de forma autónoma, gestionando desde modelos de NLP hasta infraestructura en la nube. \\
Mi investigación en el National Physical Laboratory produjo un novedoso indicador de escala para la detección de puntos de inflexión en series temporales multidimensionales, publicado en \href{https://doi.org/10.1088/1748-9326/ac526f}{Environmental Research Letters (2022)}.\\
Busco un puesto remoto de ciencia de datos donde pueda construir sistemas en producción y continuar resolviendo problemas técnicos complejos.

% \vspace{.05in}

% I have  experience lecturing upto Masters level, presenting at international conferences
% and \href{https://www.linkedin.com/in/joshuaprettyman/details/publications/}{publishing in respected journals}, so I'm used
% to communicating technical concepts to a variety of audiences.

\vspace{-6pt}
\begin{center}
		\rule{0.5\columnwidth}{1pt}
\end{center}
\vspace{-10pt}

% \section{Tech Stack}
\section{Habilidades Técnicas}
\vspace{-6pt}
\begin{innerlist}
	% \item Deployment to remote UNIX server using Docker and GitHub.
	\item Uso regular (10+ años) de librerías de Python: ScikitLearn, Numpy, Pandas, NLTK, Matplotlib, Plotly.
	\item Desarrollo full stack con Bash, Python, Google Cloud Platform, SQL y JavaScript.
	\item Transformers de Huggingface y LLMs, Ollama, OpenAI Assistant API.
	\item Familiarizado con Git, GitHub, gestión de paquetes, ETL \& ELT, CI/CD, async, REST, GraphQL
	\item Experiencia con C++, MATLAB, Retool, ingeniería de datos, dashboards, computación científica.
\end{innerlist}

\vspace{-6pt}
\begin{center}
		\rule{0.5\columnwidth}{1pt}
\end{center}
\vspace{-10pt}

\section{Experiencia Profesional}

\worksubsection{Científico de Datos}{\href{https://www.blinkseo.co.uk}{Blink SEO} / \href{https://www.macaroni.works}{Macaroni Software}}{nov. 2021 a dic. 2024}

Construí una plataforma ML full-stack desde cero, aumentando la productividad del equipo en $20\times$ y permitiendo que campañas de SEO de un año se \href{https://www.linkedin.com/posts/sam-wright-17b6ab6_shopify-seo-activity-7170336529146441729-JGDn?utm_source=combined_share_message&utm_medium=member_desktop}{entregaran en semanas}.\\

\vspace{-6pt}

\begin{innerlist}
	\item Desarrollé agrupación de palabras clave mediante NLP, combinada con datos cuantitativos, reduciendo el trabajo manual de días a minutos.
	\item Arquitecturé un backend en Python (\textit{GCP Compute Engine}) integrado con \textit{BigQuery} para procesar más de 50M de puntos de datos diarios desde APIs de servicios externos.
	\item Trabajé con el equipo para replicar y luego mejorar sus procesos (basados en hojas de cálculo) mediante retroalimentación continua y despliegue.
	\item Creé un sistema de cola de trabajos (\textit{PostgreSQL}) para ejecutar tareas de incorporación de clientes, importación de datos, análisis de datos y ML de forma asíncrona en el backend.
	\item Construí dashboards interactivos (\textit{Plotly} + \textit{JavaScript}) permitiendo a los clientes identificar y priorizar oportunidades de ingresos en minutos.
	\item Integré LLMs (\textit{huggingface}) para generar sugerencias listas para usar sobre brechas de contenido, permitiendo a los clientes ver el impacto inmediato.
	\item Entregué análisis ad-hoc e ingeniería de datos para presentaciones a inversores de nivel directivo e informes internos.
\end{innerlist}

\vspace{14pt}
\worksubsection{Investigador en Ciencia de Datos}{\href{http://www.npl.co.uk}{National Physical Laboratory (NPL)}}{sep. 2015 a feb. 2021}

Produje un novedoso indicador de escala para la detección de puntos de inflexión en series temporales multidimensionales.\\

\vspace{-6pt}

\begin{innerlist}
	\item Presenté resultados en reuniones con partes interesadas en NPL, en artículos publicados y en conferencias internacionales.
	\item Desarrollé métodos para detectar puntos de inflexión en sistemas dinámicos, trabajando con grandes conjuntos de datos de series temporales y construyendo modelos estocásticos para probar hipótesis.
	\item Código escrito en \textit{MATLAB} y \textit{Python}. Visualizaciones producidas usando \textit{Matplotlib} y \textit{XMGrace}.
	\item Acceso, limpieza, interpolación e interpretación correcta de grandes conjuntos de datos meteorológicos en bruto.
\end{innerlist}


\pagebreak

\sectioncont{Experiencia Profesional}

\worksubsection{Profesor Asociado}{\href{https://www.shu.ac.uk/about-us/academic-departments/engineering-and-mathematics}{Sheffield Hallam University}}{sep. 2017 a jul. 2019}

Impartí cursos de matemáticas, estadística e informática desde Grado Básico hasta Máster.\\

\vspace{-6pt}

\begin{innerlist}
	\item Impartí un curso sobre \textit{Microsoft Excel para Negocios}.
	\item Las responsabilidades incluían preparación de lecciones, planificación de cursos, corrección y dirección de tutorías.
	\item Atendí el mostrador de ``Ayuda de Matemáticas'' tres veces por semana, proporcionando apoyo a estudiantes de todos los departamentos.
\end{innerlist}
\vspace{10pt}

\worksubsection{Desarrollador de informática (prácticas)}{\href{https://www.metoffice.gov.uk/research/foundation/informatics-lab/what-we-do}{UK Met Office Informatics Lab}}{jun. a ago. 2017}

Un proyecto de tres meses evaluando la viabilidad de usar datos en tiempo real de meteorólogos aficionados para aumentar los datos de observación oficiales. \\

\vspace{-6pt}

\begin{innerlist}
	\item Accedí a datos internos y externos mediante APIs usando Java.
	\item Limpieza de datos y análisis en Python.
	\item Trabajé como parte de un equipo y participé en reuniones diarias (daily stand-ups).
\end{innerlist}



%\vspace{.1in}
%\textbf{Informatics developer} — \href{https://www.metoffice.gov.uk/research/foundation/informatics-lab/what-we-do}{UK Met Office Informatics Lab} \hfill \emph{2017}
%\begin{innerlist}
%	\item 3 Month project involved consolidating and testing observation data.
%	\item Worked as part of a team: daily stand-ups, using Slack to coordinate.
%	\item Working in Python, accessing internal and external data via APIs.
%\end{innerlist}

%\textbf{Teaching Assistant} — University of Reading \hfill \emph{2014 - 2015}
%\vspace{.05in}
%	\begin{innerlist}
%		\item Tutorial support and teaching cover on undergraduate courses whilst undertaking my own Masters studies.
%		\item Responsibilities such as marking and exam invigilator.
%	\end{innerlist}



\vspace{1pt}
\begin{center}
		\rule{0.5\columnwidth}{1pt}
\end{center}
\vspace{-10pt}

% \pagebreak


\section{Formación Académica}

\edusubsection{Doctorado en Matemáticas}{University of Reading}{}
\hspace{2pt} \textbf{Tesis:} \href{http://centaur.reading.ac.uk/98364/1/23022044_Prettyman_Thesis_Joshua\%20Prettyman.pdf}{\emph{Tipping Points and Early Warning Signals with Applications to Geophysical Data}}.
\vspace{4pt}
\begin{innerlist}
	\item La tesis desarrolla métodos de análisis espectral para detectar correlaciones en series temporales multidimensionales, que se utilizan como señales de alerta temprana para eventos de inflexión en sistemas dinámicos estocásticos.
	\item Publicación de tres artículos en revistas de prestigio:
		\begin{innerlist2}
			\item \textit{A novel scaling indicator of early warning signals} \href{https://doi.org/10.1209/0295-5075/121/10002}{(EPL, 2018)}
			\item \textit{Generalized early warning signals in multivariate and gridded data} \href{https://doi.org/10.1063/1.5093495}{(Chaos, 2019)}
			\item \textit{Power spectrum scaling as a measure of critical slowing down} \href{https://doi.org/10.1088/1748-9326/ac526f}{(ERL, 2022)}
		\end{innerlist2}
	\item Implementé métodos derivados matemáticamente de forma numérica mediante discretización y linealización.
	\item Participación en varias conferencias y talleres internacionales.
	\item Prácticas de tres meses enfocadas en datos en la UK Met Office (ver sección \textit{Experiencia Profesional}).
\end{innerlist}
\vspace{14pt}

\edusubsection{M.Res. Matemáticas}{Imperial College London}{Distinción}
\hspace{2pt} \textbf{Tesis:} \href{https://github.com/DrPrettyman/dissertations/blob/main/MresDiss_Prettyman.pdf}{\emph{Mesh Generation Using Solutions of an Optimal Transport Problem}}.
\vspace{4pt}
\begin{innerlist}
	% \item Dissertation title: \href{https://github.com/DrPrettyman/dissertations/blob/main/MresDiss_Prettyman.pdf}{\emph{Adaptive Mesh Generation}}.
	\item La tesis utiliza soluciones a una linealización de la ecuación de Monge-Ampère para generar una malla adaptativa y distribuida de forma óptima para modelos numéricos.
	\item Trabajo en C++; monitorización del uso de CPU y tiempo de convergencia en varios escenarios.
	\item Participación en talleres semanales de habilidades blandas y presentaciones internas mensuales.
	\item Asistencia a seminarios semanales de Matemáticas del Planeta Tierra organizados por el Centro de Formación Doctoral.
	\item Completé cursos impartidos en Python, R, Sistemas Dinámicos, Ecuaciones diferenciales estocásticas, Análisis Bayesiano y Probabilidad.
\end{innerlist}
\vspace{14pt}

\edusubsection{MA Matemáticas}{University of Edinburgh}{Primera Clase con Honores}
\hspace{2pt} \textbf{Tesis:} \href{https://github.com/DrPrettyman/dissertations/blob/main/MaDiss_Prettyman.pdf}{\emph{Growth and its Applications in Graph Theory}}.
\vspace{4pt}
\begin{innerlist}
	\item La tesis explora patrones en caminos a través de dígrafos conectados.
	\item Escribí un programa con una interfaz de usuario simple para investigar las propiedades de estos dígrafos: esto se convirtió posteriormente en una herramienta de enseñanza, ganándome una carta de agradecimiento del director.
	\item Cursos obligatorios para el grado de Matemáticas Puras además de módulos opcionales en informática y educación matemática.
	\item Divulgación matemática en escuelas locales, eventos en museos y seminarios universitarios.
\end{innerlist}


\end{document}

%%%%%%%%%%%%%%%%%%%%%%%%%% End CV Document %%%%%%%%%%%%%%%%%%%%%%%%%%%%%
